%%%%%%%%%%%%%%%%%%%%%%%%%%%%%%%%%%%%%%%%%%%%%%%%%%%%%%%%%%%%%%%%%%%%%%
% Lớp 8 | Chương II. Hằng đẳng thức đáng nhớ và ứng dụng | Lập phương của một tổng. Lập phương của một hiệu
% Dạng 8D2?7-1: Khai triển lập phương của một tổng, một hiệu
%%%%%%%%%%%%%%%%%%%%%%%%%%%%%%%%%%%%%%%%%%%%%%%%%%%%%%%%%%%%%%%%%%%%%%


% nguồn: dan_tu_AI (nhập 16/11/2026)
\begin{bt}%[8D2H7-1]%[ma:591f684f]
	Khai triển biểu thức $(x+4)^3$.
	\loigiai{
		Ta có:
		\begin{align*} 
			(x+4)^3 &= x^3 + 3 \cdot x^2 \cdot 4 + 3 \cdot x \cdot 4^2 + 4^3 \\
			&= x^3 + 12x^2 + 48x + 64.
		\end{align*}
	}
\end{bt}

% nguồn: dan_tu_AI (nhập 16/11/2026)
\begin{bt}%[8D2H7-1]%[ma:9a976b32]
	Khai triển biểu thức $(2x-3)^3$.
	\loigiai{
		Ta có:
		\begin{align*} 
			(2x-3)^3 &= (2x)^3 - 3 \cdot (2x)^2 \cdot 3 + 3 \cdot (2x) \cdot 3^2 - 3^3 \\
			&= 8x^3 - 3 \cdot 4x^2 \cdot 3 + 6x \cdot 9 - 27 \\
			&= 8x^3 - 36x^2 + 54x - 27.
		\end{align*}
	}
\end{bt}

% nguồn: dan_tu_AI (nhập 16/11/2026)
\begin{bt}%[8D2H7-1]%[ma:c51bff52]
	Khai triển biểu thức $(x+3y)^3$.
	\loigiai{
		Ta có:
		\begin{align*} 
			(x+3y)^3 &= x^3 + 3 \cdot x^2 \cdot (3y) + 3 \cdot x \cdot (3y)^2 + (3y)^3 \\
			&= x^3 + 9x^2y + 3 \cdot x \cdot 9y^2 + 27y^3 \\
			&= x^3 + 9x^2y + 27xy^2 + 27y^3.
		\end{align*}
	}
\end{bt}

% nguồn: dan_tu_AI (nhập 16/11/2026)
\begin{bt}%[8D2H7-1]%[ma:84247fad]
	Khai triển biểu thức $(3x-y)^3$.
	\loigiai{
		Ta có:
		\begin{align*} 
			(3x-y)^3 &= (3x)^3 - 3 \cdot (3x)^2 \cdot y + 3 \cdot (3x) \cdot y^2 - y^3 \\
			&= 27x^3 - 3 \cdot 9x^2 \cdot y + 9xy^2 - y^3 \\
			&= 27x^3 - 27x^2y + 9xy^2 - y^3.
		\end{align*}
	}
\end{bt}

% nguồn: dan_tu_AI (nhập 16/11/2026)
\begin{bt}%[8D2H7-1]%[ma:fc66a23e]
	Khai triển biểu thức $\left(2x+\dfrac{1}{2}\right)^3$.
	\loigiai{
		Ta có:
		\begin{align*} 
			\left(2x+\dfrac{1}{2}\right)^3 &= (2x)^3 + 3 \cdot (2x)^2 \cdot \dfrac{1}{2} + 3 \cdot (2x) \cdot \left(\dfrac{1}{2}\right)^2 + \left(\dfrac{1}{2}\right)^3 \\
			&= 8x^3 + 3 \cdot 4x^2 \cdot \dfrac{1}{2} + 6x \cdot \dfrac{1}{4} + \dfrac{1}{8} \\
			&= 8x^3 + 6x^2 + \dfrac{3}{2}x + \dfrac{1}{8}.
		\end{align*}
	}
\end{bt}

% nguồn: dan_tu_AI (nhập 16/11/2026)
\begin{bt}%[8D2V7-1]%[ma:4b5fad4c]
	Một hộp quà có dạng hình lập phương với độ dài cạnh là $3x+2$ (cm). Viết đa thức khai triển biểu thị thể tích của hộp quà đó.
	\loigiai{
		Thể tích của hộp quà hình lập phương là:
		\[ V = (3x+2)^3. \]
		Áp dụng hằng đẳng thức lập phương của một tổng, ta khai triển đa thức:
		\begin{align*} 
			V &= (3x)^3 + 3 \cdot (3x)^2 \cdot 2 + 3 \cdot (3x) \cdot 2^2 + 2^3 \\
			&= 27x^3 + 3 \cdot 9x^2 \cdot 2 + 9x \cdot 4 + 8 \\
			&= 27x^3 + 54x^2 + 36x + 8 \text{ (cm}^3\text{)}.
		\end{align*}
		Vậy đa thức biểu thị thể tích hộp quà là $27x^3 + 54x^2 + 36x + 8$.
	}
\end{bt}
